\scanpage{45}
Calculons la \uline{trace} de la représentation \unsure{dévissée|divisée|ci-dessus}, on trouve pour $\varphi\deleted{\unsure{\xi|g}}\in D/I$
\[
 \operatorname{Tr}_{E^{\mathord{\text{\tikz[x=.66em,y=.68em,baseline=0pt]{\draw[line join=round] (.80,.95)--(.12,.24)--(.83,.34)--(.14,-.37);}}}}}(\varphi)\deleted{\unsure{\mathord{\text{\tikz[x=.66em,y=.68em,baseline=0pt]{\draw[line join=round] (.80,.95)--(.12,.24)--(.83,.34)--(.14,-.37);}}}|g}}
 =\sum_i\operatorname{Tr}_{E_i^{\mathord{\text{\tikz[x=.66em,y=.68em,baseline=0pt]{\draw[line join=round] (.80,.95)--(.12,.24)--(.83,.34)--(.14,-.37);}}}}}(\varphi)
\]

Or si $\deleted{\unsure{I|G}}$ $U$ opère trivialmt sur $E$, \deleted{qui est} alors un module sur $D/U$, on a
\[
 \deleted{\unsure{\text{en}|\text{on}}}\ E^{\mathord{\text{\tikz[x=.66em,y=.68em,baseline=0pt]{\draw[line join=round] (.80,.95)--(.12,.24)--(.83,.34)--(.14,-.37);}}}}
 =\operatorname{Im}\deleted{\unsure{\tfrac{1}{N}|\mathord{\text{\tikz[x=.66em,y=.68em,baseline=0pt]{\draw[line join=round] (.80,.95)--(.12,.24)--(.83,.34)--(.14,-.37);}}}}}\underset{\text{\inserted{\unsure{projection|projecteur|régulière}}}}{\pi}\deleted{\unsure{E|I}}
 \qquad \pi=\frac{1}{N}\sum_{i\in I/U}e_i \quad .\ N=\operatorname{Card} I/U
\]
donc si $g\in I/U\longrightarrow\varphi$ \unsure{on aura|on a|\ldots\ \ldots}
\[
 \operatorname{Tr}_{E^{\mathord{\text{\tikz[x=.66em,y=.68em,baseline=0pt]{\draw[line join=round] (.80,.95)--(.12,.24)--(.83,.34)--(.14,-.37);}}}}}(\varphi)=\operatorname{Tr}_E(\deleted{\unsure{I|1}}\pi g)
 \deleted{=\frac{1}{N}\sum_{\substack{g\to\varphi\\\unsure{\text{dans }I|g\in I}}}\operatorname{Tr}_E(}
 \qquad\text{d'où}
\]
\[
 \boxed{\displaystyle\operatorname{Tr}_{E^{\mathord{\text{\tikz[x=.66em,y=.68em,baseline=0pt]{\draw[line join=round] (.80,.95)--(.12,.24)--(.83,.34)--(.14,-.37);}}}}}(\varphi)=
 \frac{1}{N}\sum_{\substack{\psi_\alpha\in D/U\\\unsure{\substack{\deleted{i}\\\psi_\alpha}|\bar{\psi}_\alpha|\substack{\mathord{-}\\\psi_\alpha}|\psi_\alpha}\to\varphi\in D/I}}\operatorname{Tr}_E\psi_\alpha}
\]

Comme \inserted{les deux \unsure{mbres|membres} de} cette formule sont linéaires en $E$, ça reste valable pour tt $E$ ; il faut simplemt noter que $\operatorname{Tr}_E\psi$ est défini \deleted{\unsure{même si|si}} pour $\psi\deleted{\unsure{\alpha|i}}\in D/U$, si $U$ opère de façon unipotente, comme $\operatorname{Tr}_E(\unsure{\tau|\gamma|\psi}')$ pour $\tau'\in D$ se restreignant en $\tau$ : cette valeur ne \uline{dépend pas} du
% archive sheet number: 44
